% !TeX root = closed-systems-from-comparison-completeness.tex
% ============================================================
\chapter*{Preface}
\addcontentsline{toc}{chapter}{Preface}
\label{ch:preface}
% ============================================================
This book is a theorem-first monograph on closed relational systems.
Its objective is precise: start from intrinsic comparison data alone,
assume no background geometry or external frame, and classify exactly
which global structures are determined.

The work is not a reinterpretation of standard geometric physics in new
language. It is an admissibility test. At each stage, a structure is
admitted only if it is internally recoverable from the relational system
itself, or else named explicitly as boundary data for a later realization.
What survives that test forms the admissible core of the theory; what does
not survive is not discarded, but is reclassified as an interface condition,
a realization choice, or an open reconstruction problem.

The resulting architecture is rigid. Closure determines quotient semantics,
extension freedom collapses to two loci, transport obstruction first
stabilizes in a quadratic layer, and smooth realization turns that same
layer into curvature. In the intrinsic--smooth interface, faithfulness of
the smooth reading is decomposed into a torsion condition and one intrinsic
detectability axiom, both stated on the tower itself, and is discharged on a
canonical realization built from the tower. The later geometric and
field-theoretic chapters therefore appear as downstream consequences of
the same intrinsic closure mechanism.

The manuscript should be read in this order. First, comparison data and
closed-world admissibility are fixed. Second, the semantic quotient and the
two possible enrichment loci are classified. Third, the morphism-level
transport residue is isolated and shown to have a first stable quadratic
carrier. Fourth, smooth geometry is introduced only as a realized reading of
that carrier, under explicitly stated interface conditions. Finally, the
macroscopic, electromagnetic, scalar, matter-sector, and state-side
reconstruction statements are placed according to the role they play in that
chain.

The manuscript is organized as a dependency chain. Early chapters prove
compactness and semantic closure; middle chapters classify enrichment and
obstruction; later chapters derive dynamics, curvature, causal scaling,
Hilbert structure, Einstein compatibility, connection-first
reconstruction, electromagnetic quantization, scalar-sector closure, and
matter-sector realization. The final part records the resulting
state-side reconstruction problem: the raw limit observable algebra is
larger than any realized manifold reading, the excess is characterized,
and the regular sector is identified with the realized algebra under a
named scale-coherence condition. The appendix then records quantitative
boundary data and a concrete confinement realization of the scalar
channel fixed in the main stack.

The theorem statements, definitions, and scoped hypotheses are the binding
claims of the text. Expository summaries are included only to make the
dependency chain readable. When a result depends on a smooth realization,
a ghost-free regime, a scale-coherence condition, or an observer-side
semantic readout, that dependence is stated as part of the result rather
than absorbed into the background.

For readers who want the exact dependency map before technical details,
Chapter~0 provides the complete theorem-level dependency sequence.
